\scanpage{16}
\mspage{7}
% archive sheet number: 16

\uline{\unsure{Mog|Mag} $L$ 2}\quad\unsure{$\checkmark$}\quad Soient $X,Y\in M$, \deleted{tels deux} tels que
\[
\forall\,x\in\operatorname{omb}(X)^\circ,y\in\operatorname{omb}(Y)^\circ,\ \text{on a}\quad x\parallel y,
\]
Alors $X\mathrel{|\unsure{\circ|0}|}Y$.

\margin{implique Mog $L$ 1 (à cause de l'antiréflexivité de $\mathrel{|\circ|}$)}
% slanted

Cela signifie donc que $\mathrel{|\circ|}$ sur $M$ se déduit de la \deleted{\unsure{construction}} relation $\parallel$ sur $\unsure{\mathscr L|L}$. Il semble que \inserted{dans les cas pratiques,} chaque fois qu'on a Mog $L$ 1, on a aussi Mog $L$ 2, qui semble donc l'axiome le plus utile.

\uline{Proposition}-\uline{\unsure{Scholie|Scolie}}. Se donner \inserted{sur $M$} \unsure{une|un} \inserted{structure de} magma ``strictement divisible'' i.e. satisfaisant : Mog $L$ 2, revient au même que de se donner $\leq$, $\ll$ satisfaisant Mog 1, Mog 2 p.1, \inserted{la condition Mog $L$ 1 (p.6)} \uline{plus} une relation
\[
x\parallel y\qquad\text{relation \uline{sur $\mathscr L$}}
\]
satisfaisant les seules conditions

\uline{Mog div 3}\quad La relation $x\parallel y$ sur $\mathscr L$ est symétrique et antiréflexive.

\uline{Mog div 4}\quad Si $Y,Z\leq X$ dans $M$, $Y\neq Z$, et si $y\in\operatorname{omb}(Y)^\circ$, $z\in\operatorname{omb}(Z)^\circ$, alors $y\parallel z$.

En effet, dans ce cas on \uline{définit} $\mathrel{|\circ|}$ dans $M$ par la condition dans Mog $L$ 2, et on vérifie immédiatement que \deleted{elle} celle-ci satisfait Mog $3^*$, Mog 4, et même

\margin{mais pour l'antiréflexivité il faut \unsure{supposer} les $\operatorname{omb}(\unsure{X_i|X})^\circ\neq\varnothing$}
% slanted
